% ============================================================
% Punkt 2: Opis techniczny projektu
% Projekt układu sumatora pełnego (Full Adder)
% z minimalizacją bramek logicznych
% ------------------------------------------------------------
% Plik przeznaczony do włączenia w dokumencie głównym Overleaf:
%   % ============================================================
% Punkt 2: Opis techniczny projektu
% Projekt układu sumatora pełnego (Full Adder)
% z minimalizacją bramek logicznych
% ------------------------------------------------------------
% Plik przeznaczony do włączenia w dokumencie głównym Overleaf:
%   % ============================================================
% Punkt 2: Opis techniczny projektu
% Projekt układu sumatora pełnego (Full Adder)
% z minimalizacją bramek logicznych
% ------------------------------------------------------------
% Plik przeznaczony do włączenia w dokumencie głównym Overleaf:
%   % ============================================================
% Punkt 2: Opis techniczny projektu
% Projekt układu sumatora pełnego (Full Adder)
% z minimalizacją bramek logicznych
% ------------------------------------------------------------
% Plik przeznaczony do włączenia w dokumencie głównym Overleaf:
%   \input{02_opis_techniczny}
% Wymagane pakiety w dokumencie głównym:
%   \usepackage[utf8]{inputenc}  \usepackage[T1]{polski}
%   \usepackage{amsmath, amssymb, array, booktabs}
% ============================================================

\section{Opis techniczny projektu}

\subsection{Założenia projektowe}

Celem projektu jest zaprojektowanie układu kombinacyjnego realizującego funkcję sumatora pełnego, a następnie zminimalizowanie liczby wykorzystanych bramek logicznych.

Projektowany układ będzie dodawał trzy bity wejściowe: dwa bity argumentów oraz bit przeniesienia z niższego rzędu:
\begin{equation*}
a, \qquad b, \qquad c_{in},
\end{equation*}
gdzie $c_{in}$ oznacza przeniesienie generowane przez poprzedni (młodszy) stopień kaskady.

Układ posiada trzy wejścia binarne oraz dwa wyjścia. Wyjścia będą informowały o wartości bitu sumy oraz o przeniesieniu do rzędu wyższego.

Przyjęto następujące założenia:
\begin{itemize}
    \item układ dodaje trzy bity wejściowe: $a$, $b$ oraz $c_{in}$,
    \item wszystkie wejścia przyjmują wartości logiczne 0 lub 1,
    \item układ jest układem kombinacyjnym,
    \item układ posiada dwa wyjścia: bit sumy $s$ oraz przeniesienie $c_{out}$,
    \item stan wyjścia $s = 1$ oznacza, że suma wejść jest nieparzysta,
    \item stan wyjścia $c_{out} = 1$ oznacza, że $a + b + c_{in} \geq 2$,
    \item podstawowym celem projektu jest minimalizacja liczby bramek logicznych,
    \item poprawność działania zostanie sprawdzona za pomocą symulacji.
\end{itemize}

Minimalizacja układu ma na celu zmniejszenie jego złożoności. Mniejsza liczba zastosowanych bramek oznacza prostszy schemat logiczny, mniejsze zapotrzebowanie na elementy oraz potencjalnie mniejsze opóźnienia propagacji sygnału.

\subsection{Specyfikacja wejść i wyjść}

Projektowany sumator pełny posiada trzy wejścia:
\begin{itemize}
    \item $a$ -- pierwszy bit argumentu,
    \item $b$ -- drugi bit argumentu,
    \item $c_{in}$ -- przeniesienie z niższego rzędu.
\end{itemize}

Wyjściami układu są:
\begin{itemize}
    \item $s$ -- bit sumy,
    \item $c_{out}$ -- przeniesienie do rzędu wyższego.
\end{itemize}

Zależność pomiędzy wejściami i wyjściami można zapisać jako:
\begin{equation*}
s, \, c_{out} = f(a, b, c_{in}).
\end{equation*}

Przykładowo, dla:
\begin{equation*}
a = 1, \qquad b = 1, \qquad c_{in} = 0,
\end{equation*}
otrzymujemy:
\begin{equation*}
s = 0, \qquad c_{out} = 1.
\end{equation*}

Natomiast dla:
\begin{equation*}
a = 1, \qquad b = 1, \qquad c_{in} = 1,
\end{equation*}
otrzymujemy:
\begin{equation*}
s = 1, \qquad c_{out} = 1.
\end{equation*}

\vspace{0.5em}
\noindent Tabela sygnałów:

\begin{center}
\begin{tabular}{|c|c|l|}
\hline
\textbf{Sygnał} & \textbf{Rodzaj} & \textbf{Opis} \\
\hline
$a$       & wejście & pierwszy bit argumentu \\
$b$       & wejście & drugi bit argumentu \\
$c_{in}$  & wejście & przeniesienie z niższego rzędu \\
$s$       & wyjście & bit sumy \\
$c_{out}$ & wyjście & przeniesienie do rzędu wyższego \\
\hline
\end{tabular}
\end{center}

\subsection{Analiza logiczna układu}

Sumator pełny wykonuje arytmetyczne dodawanie trzech bitów, w wyniku czego może powstać wartość z zakresu od 0 do 3, wymagająca dwóch bitów zapisu: $c_{out}s$.

Aby suma była równa 1, nieparzysta liczba sygnałów wejściowych musi przyjmować wartość wysoką. Warunek ten realizuje dwuargumentowa operacja alternatywy rozłącznej, zastosowana dwukrotnie:
\begin{equation*}
s = a \oplus b \oplus c_{in}.
\end{equation*}

Przeniesienie pojawia się wtedy, gdy co najmniej dwa z trzech wejść znajdują się w stanie wysokim. Funkcję wyjściową przeniesienia można zapisać jako:
\begin{equation*}
c_{out} = a \cdot b \vee a \cdot c_{in} \vee b \cdot c_{in}.
\end{equation*}

Jest to tzw. funkcja większości, która przyjmuje wartość 1, gdy większość (co najmniej dwa) sygnałów wejściowych ma wartość 1.

Wprowadzając oznaczenie pomocnicze:
\begin{equation*}
c_1 = a \oplus b,
\end{equation*}
można obydwie funkcje zapisać w postaci:
\begin{equation*}
s = c_1 \oplus c_{in}, \qquad c_{out} = a \cdot b \vee c_1 \cdot c_{in}.
\end{equation*}

Jeżeli nie korzystamy z bramek XOR, alternatywę rozłączną można zapisać przy pomocy podstawowych operacji AND, OR i NOT:
\begin{equation*}
a \oplus b = a \cdot b' \vee a' \cdot b.
\end{equation*}

Stąd:
\begin{equation*}
s = a \cdot b' \cdot c_{in}' \vee a \cdot b \cdot c_{in} \vee a' \cdot b \cdot c_{in}' \vee a' \cdot b' \cdot c_{in}.
\end{equation*}

Ta postać będzie szczególnie istotna podczas dalszej analizy i minimalizacji układu.

\subsection{Tablica prawdy}

Ponieważ układ posiada trzy wejścia binarne, liczba możliwych kombinacji wynosi:
\begin{equation*}
2^3 = 8.
\end{equation*}

Pełna tablica prawdy przedstawia się następująco:

\begin{center}
\begin{tabular}{|c|c|c|c|c|c|}
\hline
$a$ & $b$ & $c_{in}$ & $a+b+c_{in}$ & $s$ & $c_{out}$ \\
\hline
0 & 0 & 0 & 0 & 0 & 0 \\
0 & 0 & 1 & 1 & 1 & 0 \\
0 & 1 & 0 & 1 & 1 & 0 \\
0 & 1 & 1 & 2 & 0 & 1 \\
1 & 0 & 0 & 1 & 1 & 0 \\
1 & 0 & 1 & 2 & 0 & 1 \\
1 & 1 & 0 & 2 & 0 & 1 \\
1 & 1 & 1 & 3 & 1 & 1 \\
\hline
\end{tabular}
\end{center}

Wartość wyjścia $s = 1$ występuje wyłącznie dla kombinacji:
\begin{equation*}
(0,0,1), \quad (0,1,0), \quad (1,0,0), \quad (1,1,1).
\end{equation*}

Funkcję w postaci sumy iloczynów mintermów można więc zapisać jako:
\begin{equation*}
s = \Sigma m(1, 2, 4, 7).
\end{equation*}

Po rozpisaniu:
\begin{equation*}
s = a' \cdot b' \cdot c_{in} \vee a' \cdot b \cdot c_{in}' \vee a \cdot b' \cdot c_{in}' \vee a \cdot b \cdot c_{in}.
\end{equation*}

Analogicznie dla wyjścia przeniesienia, gdzie $c_{out} = 1$ dla kombinacji $(0,1,1)$, $(1,0,1)$, $(1,1,0)$ oraz $(1,1,1)$:
\begin{equation*}
c_{out} = \Sigma m(3, 5, 6, 7),
\end{equation*}
\begin{equation*}
c_{out} = a' \cdot b \cdot c_{in} \vee a \cdot b' \cdot c_{in} \vee a \cdot b \cdot c_{in}' \vee a \cdot b \cdot c_{in}.
\end{equation*}

Jest to postać kanoniczna, która będzie punktem wyjścia do minimalizacji.

\subsection{Mapa Karnaugha}

Do uproszczenia funkcji wykorzystujemy mapy Karnaugha dla trzech zmiennych, osobno dla każdego z wyjść.

Przyjmujemy:
\begin{itemize}
    \item wiersze: pary zmiennych $ab$,
    \item kolumny: zmienna $c_{in}$,
    \item wiersze uporządkowane w kolejności kodu Graya: $00, 01, 11, 10$.
\end{itemize}

Mapa Karnaugha dla wyjścia $s$ przyjmuje postać:

\begin{center}
\begin{tabular}{|c|c|c|}
\hline
$ab \backslash c_{in}$ & \textbf{0} & \textbf{1} \\
\hline
\textbf{00} & 0 & 1 \\
\hline
\textbf{01} & 1 & 0 \\
\hline
\textbf{11} & 0 & 1 \\
\hline
\textbf{10} & 1 & 0 \\
\hline
\end{tabular}
\end{center}

Jedynki znajdują się więc na przekątnej mapy. W tym przypadku nie można połączyć sąsiednich jedynek w większe grupy, ponieważ żadna z nich nie posiada sąsiedniej jedynki zgodnie z zasadami map Karnaugha. Oznacza to, że bez wykorzystania dodatkowych zależności funkcja $s$ pozostaje w postaci czterech iloczynów trójzmiennych. Jednocześnie sama struktura problemu pozwala zauważyć, że funkcję można znacznie efektywniej zrealizować, rozpoznając w niej dwukrotne zastosowanie alternatywy rozłącznej (XOR).

Mapa Karnaugha dla wyjścia $c_{out}$ przyjmuje postać:

\begin{center}
\begin{tabular}{|c|c|c|}
\hline
$ab \backslash c_{in}$ & \textbf{0} & \textbf{1} \\
\hline
\textbf{00} & 0 & 0 \\
\hline
\textbf{01} & 0 & 1 \\
\hline
\textbf{11} & 1 & 1 \\
\hline
\textbf{10} & 0 & 1 \\
\hline
\end{tabular}
\end{center}

Dla funkcji $c_{out}$ możliwe jest utworzenie trzech grup po dwie jedynki: pary w wierszu $ab = 11$, pary w kolumnie $c_{in} = 1$ obejmującej wiersze $01$ i $11$ oraz pary w kolumnie $c_{in} = 1$ obejmującej wiersze $11$ i $10$. W wyniku grupowania otrzymujemy zredukowaną postać funkcji większości:
\begin{equation*}
c_{out} = a \cdot b \vee b \cdot c_{in} \vee a \cdot c_{in}.
\end{equation*}

\subsection{Minimalizacja funkcji logicznej}

Dla wyjścia sumy, na podstawie analizy struktury funkcji, otrzymujemy:
\begin{equation*}
s = a \oplus b \oplus c_{in}.
\end{equation*}

Wprowadzając sygnał pomocniczy:
\begin{equation*}
c_1 = a \oplus b,
\end{equation*}
zapis ten upraszcza się do:
\begin{equation*}
s = c_1 \oplus c_{in}.
\end{equation*}

Dla wyjścia przeniesienia, w wyniku grupowania na mapie Karnaugha:
\begin{equation*}
c_{out} = a \cdot b \vee a \cdot c_{in} \vee b \cdot c_{in}.
\end{equation*}

Alternatywnie, wykorzystując wspólny sygnał $c_1$, przeniesienie można zapisać jako:
\begin{equation*}
c_{out} = a \cdot b \vee c_1 \cdot c_{in}.
\end{equation*}

Jest to postać znacznie bardziej przejrzysta od postaci kanonicznej.

Jeżeli w realizacji można zastosować bramki XOR, każdy z członów może zostać zrealizowany pojedynczą bramką:
\begin{equation*}
c_1 = a \oplus b, \qquad s = c_1 \oplus c_{in}.
\end{equation*}

Ostateczna realizacja wymaga więc:
\begin{itemize}
    \item 2 $\times$ XOR ($c_1$ oraz $s$),
    \item 2 $\times$ AND2 ($a \cdot b$ oraz $c_1 \cdot c_{in}$),
    \item 1 $\times$ OR2 (sumowanie członów przeniesienia).
\end{itemize}

Czyli łącznie 5 bramek.

To rozwiązanie będzie naszym wariantem zoptymalizowanym pod względem liczby bramek, o ile w dalszej części przyjmiemy XOR jako dostępny element logiczny.

\subsection{Porównanie liczby bramek przed i po minimalizacji}

Dla realizacji funkcji w sposób bezpośredni, na podstawie postaci kanonicznej, konieczne byłoby zastosowanie wielu bramek AND, NOT oraz OR.

Postać kanoniczna wyjścia sumy:
\begin{equation*}
s = a' \cdot b' \cdot c_{in} \vee a' \cdot b \cdot c_{in}' \vee a \cdot b' \cdot c_{in}' \vee a \cdot b \cdot c_{in}
\end{equation*}
wymaga utworzenia czterech iloczynów trójzmiennych, a następnie ich zsumowania. Realizacja ta wymaga: 3 bramek NOT (negacje $a$, $b$, $c_{in}$), 4 bramek AND3 oraz 1 bramki OR4, czyli łącznie 8 bramek. Analogicznie dla wyjścia $c_{out}$ kolejne 8 bramek, przy czym bramki NOT mogą być wspólne.

Łącznie realizacja bezpośrednia wymaga około 13--16 bramek logicznych o głębokości trzech poziomów.

Po zastosowaniu zależności wynikającej z alternatywy rozłącznej oraz grupowania na mapach Karnaugha otrzymujemy:
\begin{equation*}
s = a \oplus b \oplus c_{in}, \qquad c_{out} = a \cdot b \vee c_1 \cdot c_{in}, \qquad c_1 = a \oplus b.
\end{equation*}

W rezultacie układ można zrealizować przy użyciu dwóch bramek XOR, dwóch bramek AND oraz jednej bramki OR, czyli łącznie 5 bramek.

Porównanie obu wariantów przedstawia tabela:

\begin{center}
\begin{tabular}{|l|c|c|}
\hline
\textbf{Wariant realizacji} & \textbf{Liczba bramek} & \textbf{Głębokość logiczna} \\
\hline
Postać kanoniczna (SOP) & ok. 13--16 & 3 poziomy \\
Wariant zminimalizowany & 5 & 3 poziomy \\
\hline
\end{tabular}
\end{center}

Zastosowanie wspólnego sygnału $c_1 = a \oplus b$ pozwala dodatkowo wyeliminować powtarzanie tej samej operacji dla obu wyjść, co stanowi dodatkową oszczędność w porównaniu z realizacją, w której każda funkcja minimalizowana jest niezależnie.

% Wymagane pakiety w dokumencie głównym:
%   \usepackage[utf8]{inputenc}  \usepackage[T1]{polski}
%   \usepackage{amsmath, amssymb, array, booktabs}
% ============================================================

\section{Opis techniczny projektu}

\subsection{Założenia projektowe}

Celem projektu jest zaprojektowanie układu kombinacyjnego realizującego funkcję sumatora pełnego, a następnie zminimalizowanie liczby wykorzystanych bramek logicznych.

Projektowany układ będzie dodawał trzy bity wejściowe: dwa bity argumentów oraz bit przeniesienia z niższego rzędu:
\begin{equation*}
a, \qquad b, \qquad c_{in},
\end{equation*}
gdzie $c_{in}$ oznacza przeniesienie generowane przez poprzedni (młodszy) stopień kaskady.

Układ posiada trzy wejścia binarne oraz dwa wyjścia. Wyjścia będą informowały o wartości bitu sumy oraz o przeniesieniu do rzędu wyższego.

Przyjęto następujące założenia:
\begin{itemize}
    \item układ dodaje trzy bity wejściowe: $a$, $b$ oraz $c_{in}$,
    \item wszystkie wejścia przyjmują wartości logiczne 0 lub 1,
    \item układ jest układem kombinacyjnym,
    \item układ posiada dwa wyjścia: bit sumy $s$ oraz przeniesienie $c_{out}$,
    \item stan wyjścia $s = 1$ oznacza, że suma wejść jest nieparzysta,
    \item stan wyjścia $c_{out} = 1$ oznacza, że $a + b + c_{in} \geq 2$,
    \item podstawowym celem projektu jest minimalizacja liczby bramek logicznych,
    \item poprawność działania zostanie sprawdzona za pomocą symulacji.
\end{itemize}

Minimalizacja układu ma na celu zmniejszenie jego złożoności. Mniejsza liczba zastosowanych bramek oznacza prostszy schemat logiczny, mniejsze zapotrzebowanie na elementy oraz potencjalnie mniejsze opóźnienia propagacji sygnału.

\subsection{Specyfikacja wejść i wyjść}

Projektowany sumator pełny posiada trzy wejścia:
\begin{itemize}
    \item $a$ -- pierwszy bit argumentu,
    \item $b$ -- drugi bit argumentu,
    \item $c_{in}$ -- przeniesienie z niższego rzędu.
\end{itemize}

Wyjściami układu są:
\begin{itemize}
    \item $s$ -- bit sumy,
    \item $c_{out}$ -- przeniesienie do rzędu wyższego.
\end{itemize}

Zależność pomiędzy wejściami i wyjściami można zapisać jako:
\begin{equation*}
s, \, c_{out} = f(a, b, c_{in}).
\end{equation*}

Przykładowo, dla:
\begin{equation*}
a = 1, \qquad b = 1, \qquad c_{in} = 0,
\end{equation*}
otrzymujemy:
\begin{equation*}
s = 0, \qquad c_{out} = 1.
\end{equation*}

Natomiast dla:
\begin{equation*}
a = 1, \qquad b = 1, \qquad c_{in} = 1,
\end{equation*}
otrzymujemy:
\begin{equation*}
s = 1, \qquad c_{out} = 1.
\end{equation*}

\vspace{0.5em}
\noindent Tabela sygnałów:

\begin{center}
\begin{tabular}{|c|c|l|}
\hline
\textbf{Sygnał} & \textbf{Rodzaj} & \textbf{Opis} \\
\hline
$a$       & wejście & pierwszy bit argumentu \\
$b$       & wejście & drugi bit argumentu \\
$c_{in}$  & wejście & przeniesienie z niższego rzędu \\
$s$       & wyjście & bit sumy \\
$c_{out}$ & wyjście & przeniesienie do rzędu wyższego \\
\hline
\end{tabular}
\end{center}

\subsection{Analiza logiczna układu}

Sumator pełny wykonuje arytmetyczne dodawanie trzech bitów, w wyniku czego może powstać wartość z zakresu od 0 do 3, wymagająca dwóch bitów zapisu: $c_{out}s$.

Aby suma była równa 1, nieparzysta liczba sygnałów wejściowych musi przyjmować wartość wysoką. Warunek ten realizuje dwuargumentowa operacja alternatywy rozłącznej, zastosowana dwukrotnie:
\begin{equation*}
s = a \oplus b \oplus c_{in}.
\end{equation*}

Przeniesienie pojawia się wtedy, gdy co najmniej dwa z trzech wejść znajdują się w stanie wysokim. Funkcję wyjściową przeniesienia można zapisać jako:
\begin{equation*}
c_{out} = a \cdot b \vee a \cdot c_{in} \vee b \cdot c_{in}.
\end{equation*}

Jest to tzw. funkcja większości, która przyjmuje wartość 1, gdy większość (co najmniej dwa) sygnałów wejściowych ma wartość 1.

Wprowadzając oznaczenie pomocnicze:
\begin{equation*}
c_1 = a \oplus b,
\end{equation*}
można obydwie funkcje zapisać w postaci:
\begin{equation*}
s = c_1 \oplus c_{in}, \qquad c_{out} = a \cdot b \vee c_1 \cdot c_{in}.
\end{equation*}

Jeżeli nie korzystamy z bramek XOR, alternatywę rozłączną można zapisać przy pomocy podstawowych operacji AND, OR i NOT:
\begin{equation*}
a \oplus b = a \cdot b' \vee a' \cdot b.
\end{equation*}

Stąd:
\begin{equation*}
s = a \cdot b' \cdot c_{in}' \vee a \cdot b \cdot c_{in} \vee a' \cdot b \cdot c_{in}' \vee a' \cdot b' \cdot c_{in}.
\end{equation*}

Ta postać będzie szczególnie istotna podczas dalszej analizy i minimalizacji układu.

\subsection{Tablica prawdy}

Ponieważ układ posiada trzy wejścia binarne, liczba możliwych kombinacji wynosi:
\begin{equation*}
2^3 = 8.
\end{equation*}

Pełna tablica prawdy przedstawia się następująco:

\begin{center}
\begin{tabular}{|c|c|c|c|c|c|}
\hline
$a$ & $b$ & $c_{in}$ & $a+b+c_{in}$ & $s$ & $c_{out}$ \\
\hline
0 & 0 & 0 & 0 & 0 & 0 \\
0 & 0 & 1 & 1 & 1 & 0 \\
0 & 1 & 0 & 1 & 1 & 0 \\
0 & 1 & 1 & 2 & 0 & 1 \\
1 & 0 & 0 & 1 & 1 & 0 \\
1 & 0 & 1 & 2 & 0 & 1 \\
1 & 1 & 0 & 2 & 0 & 1 \\
1 & 1 & 1 & 3 & 1 & 1 \\
\hline
\end{tabular}
\end{center}

Wartość wyjścia $s = 1$ występuje wyłącznie dla kombinacji:
\begin{equation*}
(0,0,1), \quad (0,1,0), \quad (1,0,0), \quad (1,1,1).
\end{equation*}

Funkcję w postaci sumy iloczynów mintermów można więc zapisać jako:
\begin{equation*}
s = \Sigma m(1, 2, 4, 7).
\end{equation*}

Po rozpisaniu:
\begin{equation*}
s = a' \cdot b' \cdot c_{in} \vee a' \cdot b \cdot c_{in}' \vee a \cdot b' \cdot c_{in}' \vee a \cdot b \cdot c_{in}.
\end{equation*}

Analogicznie dla wyjścia przeniesienia, gdzie $c_{out} = 1$ dla kombinacji $(0,1,1)$, $(1,0,1)$, $(1,1,0)$ oraz $(1,1,1)$:
\begin{equation*}
c_{out} = \Sigma m(3, 5, 6, 7),
\end{equation*}
\begin{equation*}
c_{out} = a' \cdot b \cdot c_{in} \vee a \cdot b' \cdot c_{in} \vee a \cdot b \cdot c_{in}' \vee a \cdot b \cdot c_{in}.
\end{equation*}

Jest to postać kanoniczna, która będzie punktem wyjścia do minimalizacji.

\subsection{Mapa Karnaugha}

Do uproszczenia funkcji wykorzystujemy mapy Karnaugha dla trzech zmiennych, osobno dla każdego z wyjść.

Przyjmujemy:
\begin{itemize}
    \item wiersze: pary zmiennych $ab$,
    \item kolumny: zmienna $c_{in}$,
    \item wiersze uporządkowane w kolejności kodu Graya: $00, 01, 11, 10$.
\end{itemize}

Mapa Karnaugha dla wyjścia $s$ przyjmuje postać:

\begin{center}
\begin{tabular}{|c|c|c|}
\hline
$ab \backslash c_{in}$ & \textbf{0} & \textbf{1} \\
\hline
\textbf{00} & 0 & 1 \\
\hline
\textbf{01} & 1 & 0 \\
\hline
\textbf{11} & 0 & 1 \\
\hline
\textbf{10} & 1 & 0 \\
\hline
\end{tabular}
\end{center}

Jedynki znajdują się więc na przekątnej mapy. W tym przypadku nie można połączyć sąsiednich jedynek w większe grupy, ponieważ żadna z nich nie posiada sąsiedniej jedynki zgodnie z zasadami map Karnaugha. Oznacza to, że bez wykorzystania dodatkowych zależności funkcja $s$ pozostaje w postaci czterech iloczynów trójzmiennych. Jednocześnie sama struktura problemu pozwala zauważyć, że funkcję można znacznie efektywniej zrealizować, rozpoznając w niej dwukrotne zastosowanie alternatywy rozłącznej (XOR).

Mapa Karnaugha dla wyjścia $c_{out}$ przyjmuje postać:

\begin{center}
\begin{tabular}{|c|c|c|}
\hline
$ab \backslash c_{in}$ & \textbf{0} & \textbf{1} \\
\hline
\textbf{00} & 0 & 0 \\
\hline
\textbf{01} & 0 & 1 \\
\hline
\textbf{11} & 1 & 1 \\
\hline
\textbf{10} & 0 & 1 \\
\hline
\end{tabular}
\end{center}

Dla funkcji $c_{out}$ możliwe jest utworzenie trzech grup po dwie jedynki: pary w wierszu $ab = 11$, pary w kolumnie $c_{in} = 1$ obejmującej wiersze $01$ i $11$ oraz pary w kolumnie $c_{in} = 1$ obejmującej wiersze $11$ i $10$. W wyniku grupowania otrzymujemy zredukowaną postać funkcji większości:
\begin{equation*}
c_{out} = a \cdot b \vee b \cdot c_{in} \vee a \cdot c_{in}.
\end{equation*}

\subsection{Minimalizacja funkcji logicznej}

Dla wyjścia sumy, na podstawie analizy struktury funkcji, otrzymujemy:
\begin{equation*}
s = a \oplus b \oplus c_{in}.
\end{equation*}

Wprowadzając sygnał pomocniczy:
\begin{equation*}
c_1 = a \oplus b,
\end{equation*}
zapis ten upraszcza się do:
\begin{equation*}
s = c_1 \oplus c_{in}.
\end{equation*}

Dla wyjścia przeniesienia, w wyniku grupowania na mapie Karnaugha:
\begin{equation*}
c_{out} = a \cdot b \vee a \cdot c_{in} \vee b \cdot c_{in}.
\end{equation*}

Alternatywnie, wykorzystując wspólny sygnał $c_1$, przeniesienie można zapisać jako:
\begin{equation*}
c_{out} = a \cdot b \vee c_1 \cdot c_{in}.
\end{equation*}

Jest to postać znacznie bardziej przejrzysta od postaci kanonicznej.

Jeżeli w realizacji można zastosować bramki XOR, każdy z członów może zostać zrealizowany pojedynczą bramką:
\begin{equation*}
c_1 = a \oplus b, \qquad s = c_1 \oplus c_{in}.
\end{equation*}

Ostateczna realizacja wymaga więc:
\begin{itemize}
    \item 2 $\times$ XOR ($c_1$ oraz $s$),
    \item 2 $\times$ AND2 ($a \cdot b$ oraz $c_1 \cdot c_{in}$),
    \item 1 $\times$ OR2 (sumowanie członów przeniesienia).
\end{itemize}

Czyli łącznie 5 bramek.

To rozwiązanie będzie naszym wariantem zoptymalizowanym pod względem liczby bramek, o ile w dalszej części przyjmiemy XOR jako dostępny element logiczny.

\subsection{Porównanie liczby bramek przed i po minimalizacji}

Dla realizacji funkcji w sposób bezpośredni, na podstawie postaci kanonicznej, konieczne byłoby zastosowanie wielu bramek AND, NOT oraz OR.

Postać kanoniczna wyjścia sumy:
\begin{equation*}
s = a' \cdot b' \cdot c_{in} \vee a' \cdot b \cdot c_{in}' \vee a \cdot b' \cdot c_{in}' \vee a \cdot b \cdot c_{in}
\end{equation*}
wymaga utworzenia czterech iloczynów trójzmiennych, a następnie ich zsumowania. Realizacja ta wymaga: 3 bramek NOT (negacje $a$, $b$, $c_{in}$), 4 bramek AND3 oraz 1 bramki OR4, czyli łącznie 8 bramek. Analogicznie dla wyjścia $c_{out}$ kolejne 8 bramek, przy czym bramki NOT mogą być wspólne.

Łącznie realizacja bezpośrednia wymaga około 13--16 bramek logicznych o głębokości trzech poziomów.

Po zastosowaniu zależności wynikającej z alternatywy rozłącznej oraz grupowania na mapach Karnaugha otrzymujemy:
\begin{equation*}
s = a \oplus b \oplus c_{in}, \qquad c_{out} = a \cdot b \vee c_1 \cdot c_{in}, \qquad c_1 = a \oplus b.
\end{equation*}

W rezultacie układ można zrealizować przy użyciu dwóch bramek XOR, dwóch bramek AND oraz jednej bramki OR, czyli łącznie 5 bramek.

Porównanie obu wariantów przedstawia tabela:

\begin{center}
\begin{tabular}{|l|c|c|}
\hline
\textbf{Wariant realizacji} & \textbf{Liczba bramek} & \textbf{Głębokość logiczna} \\
\hline
Postać kanoniczna (SOP) & ok. 13--16 & 3 poziomy \\
Wariant zminimalizowany & 5 & 3 poziomy \\
\hline
\end{tabular}
\end{center}

Zastosowanie wspólnego sygnału $c_1 = a \oplus b$ pozwala dodatkowo wyeliminować powtarzanie tej samej operacji dla obu wyjść, co stanowi dodatkową oszczędność w porównaniu z realizacją, w której każda funkcja minimalizowana jest niezależnie.

% Wymagane pakiety w dokumencie głównym:
%   \usepackage[utf8]{inputenc}  \usepackage[T1]{polski}
%   \usepackage{amsmath, amssymb, array, booktabs}
% ============================================================

\section{Opis techniczny projektu}

\subsection{Założenia projektowe}

Celem projektu jest zaprojektowanie układu kombinacyjnego realizującego funkcję sumatora pełnego, a następnie zminimalizowanie liczby wykorzystanych bramek logicznych.

Projektowany układ będzie dodawał trzy bity wejściowe: dwa bity argumentów oraz bit przeniesienia z niższego rzędu:
\begin{equation*}
a, \qquad b, \qquad c_{in},
\end{equation*}
gdzie $c_{in}$ oznacza przeniesienie generowane przez poprzedni (młodszy) stopień kaskady.

Układ posiada trzy wejścia binarne oraz dwa wyjścia. Wyjścia będą informowały o wartości bitu sumy oraz o przeniesieniu do rzędu wyższego.

Przyjęto następujące założenia:
\begin{itemize}
    \item układ dodaje trzy bity wejściowe: $a$, $b$ oraz $c_{in}$,
    \item wszystkie wejścia przyjmują wartości logiczne 0 lub 1,
    \item układ jest układem kombinacyjnym,
    \item układ posiada dwa wyjścia: bit sumy $s$ oraz przeniesienie $c_{out}$,
    \item stan wyjścia $s = 1$ oznacza, że suma wejść jest nieparzysta,
    \item stan wyjścia $c_{out} = 1$ oznacza, że $a + b + c_{in} \geq 2$,
    \item podstawowym celem projektu jest minimalizacja liczby bramek logicznych,
    \item poprawność działania zostanie sprawdzona za pomocą symulacji.
\end{itemize}

Minimalizacja układu ma na celu zmniejszenie jego złożoności. Mniejsza liczba zastosowanych bramek oznacza prostszy schemat logiczny, mniejsze zapotrzebowanie na elementy oraz potencjalnie mniejsze opóźnienia propagacji sygnału.

\subsection{Specyfikacja wejść i wyjść}

Projektowany sumator pełny posiada trzy wejścia:
\begin{itemize}
    \item $a$ -- pierwszy bit argumentu,
    \item $b$ -- drugi bit argumentu,
    \item $c_{in}$ -- przeniesienie z niższego rzędu.
\end{itemize}

Wyjściami układu są:
\begin{itemize}
    \item $s$ -- bit sumy,
    \item $c_{out}$ -- przeniesienie do rzędu wyższego.
\end{itemize}

Zależność pomiędzy wejściami i wyjściami można zapisać jako:
\begin{equation*}
s, \, c_{out} = f(a, b, c_{in}).
\end{equation*}

Przykładowo, dla:
\begin{equation*}
a = 1, \qquad b = 1, \qquad c_{in} = 0,
\end{equation*}
otrzymujemy:
\begin{equation*}
s = 0, \qquad c_{out} = 1.
\end{equation*}

Natomiast dla:
\begin{equation*}
a = 1, \qquad b = 1, \qquad c_{in} = 1,
\end{equation*}
otrzymujemy:
\begin{equation*}
s = 1, \qquad c_{out} = 1.
\end{equation*}

\vspace{0.5em}
\noindent Tabela sygnałów:

\begin{center}
\begin{tabular}{|c|c|l|}
\hline
\textbf{Sygnał} & \textbf{Rodzaj} & \textbf{Opis} \\
\hline
$a$       & wejście & pierwszy bit argumentu \\
$b$       & wejście & drugi bit argumentu \\
$c_{in}$  & wejście & przeniesienie z niższego rzędu \\
$s$       & wyjście & bit sumy \\
$c_{out}$ & wyjście & przeniesienie do rzędu wyższego \\
\hline
\end{tabular}
\end{center}

\subsection{Analiza logiczna układu}

Sumator pełny wykonuje arytmetyczne dodawanie trzech bitów, w wyniku czego może powstać wartość z zakresu od 0 do 3, wymagająca dwóch bitów zapisu: $c_{out}s$.

Aby suma była równa 1, nieparzysta liczba sygnałów wejściowych musi przyjmować wartość wysoką. Warunek ten realizuje dwuargumentowa operacja alternatywy rozłącznej, zastosowana dwukrotnie:
\begin{equation*}
s = a \oplus b \oplus c_{in}.
\end{equation*}

Przeniesienie pojawia się wtedy, gdy co najmniej dwa z trzech wejść znajdują się w stanie wysokim. Funkcję wyjściową przeniesienia można zapisać jako:
\begin{equation*}
c_{out} = a \cdot b \vee a \cdot c_{in} \vee b \cdot c_{in}.
\end{equation*}

Jest to tzw. funkcja większości, która przyjmuje wartość 1, gdy większość (co najmniej dwa) sygnałów wejściowych ma wartość 1.

Wprowadzając oznaczenie pomocnicze:
\begin{equation*}
c_1 = a \oplus b,
\end{equation*}
można obydwie funkcje zapisać w postaci:
\begin{equation*}
s = c_1 \oplus c_{in}, \qquad c_{out} = a \cdot b \vee c_1 \cdot c_{in}.
\end{equation*}

Jeżeli nie korzystamy z bramek XOR, alternatywę rozłączną można zapisać przy pomocy podstawowych operacji AND, OR i NOT:
\begin{equation*}
a \oplus b = a \cdot b' \vee a' \cdot b.
\end{equation*}

Stąd:
\begin{equation*}
s = a \cdot b' \cdot c_{in}' \vee a \cdot b \cdot c_{in} \vee a' \cdot b \cdot c_{in}' \vee a' \cdot b' \cdot c_{in}.
\end{equation*}

Ta postać będzie szczególnie istotna podczas dalszej analizy i minimalizacji układu.

\subsection{Tablica prawdy}

Ponieważ układ posiada trzy wejścia binarne, liczba możliwych kombinacji wynosi:
\begin{equation*}
2^3 = 8.
\end{equation*}

Pełna tablica prawdy przedstawia się następująco:

\begin{center}
\begin{tabular}{|c|c|c|c|c|c|}
\hline
$a$ & $b$ & $c_{in}$ & $a+b+c_{in}$ & $s$ & $c_{out}$ \\
\hline
0 & 0 & 0 & 0 & 0 & 0 \\
0 & 0 & 1 & 1 & 1 & 0 \\
0 & 1 & 0 & 1 & 1 & 0 \\
0 & 1 & 1 & 2 & 0 & 1 \\
1 & 0 & 0 & 1 & 1 & 0 \\
1 & 0 & 1 & 2 & 0 & 1 \\
1 & 1 & 0 & 2 & 0 & 1 \\
1 & 1 & 1 & 3 & 1 & 1 \\
\hline
\end{tabular}
\end{center}

Wartość wyjścia $s = 1$ występuje wyłącznie dla kombinacji:
\begin{equation*}
(0,0,1), \quad (0,1,0), \quad (1,0,0), \quad (1,1,1).
\end{equation*}

Funkcję w postaci sumy iloczynów mintermów można więc zapisać jako:
\begin{equation*}
s = \Sigma m(1, 2, 4, 7).
\end{equation*}

Po rozpisaniu:
\begin{equation*}
s = a' \cdot b' \cdot c_{in} \vee a' \cdot b \cdot c_{in}' \vee a \cdot b' \cdot c_{in}' \vee a \cdot b \cdot c_{in}.
\end{equation*}

Analogicznie dla wyjścia przeniesienia, gdzie $c_{out} = 1$ dla kombinacji $(0,1,1)$, $(1,0,1)$, $(1,1,0)$ oraz $(1,1,1)$:
\begin{equation*}
c_{out} = \Sigma m(3, 5, 6, 7),
\end{equation*}
\begin{equation*}
c_{out} = a' \cdot b \cdot c_{in} \vee a \cdot b' \cdot c_{in} \vee a \cdot b \cdot c_{in}' \vee a \cdot b \cdot c_{in}.
\end{equation*}

Jest to postać kanoniczna, która będzie punktem wyjścia do minimalizacji.

\subsection{Mapa Karnaugha}

Do uproszczenia funkcji wykorzystujemy mapy Karnaugha dla trzech zmiennych, osobno dla każdego z wyjść.

Przyjmujemy:
\begin{itemize}
    \item wiersze: pary zmiennych $ab$,
    \item kolumny: zmienna $c_{in}$,
    \item wiersze uporządkowane w kolejności kodu Graya: $00, 01, 11, 10$.
\end{itemize}

Mapa Karnaugha dla wyjścia $s$ przyjmuje postać:

\begin{center}
\begin{tabular}{|c|c|c|}
\hline
$ab \backslash c_{in}$ & \textbf{0} & \textbf{1} \\
\hline
\textbf{00} & 0 & 1 \\
\hline
\textbf{01} & 1 & 0 \\
\hline
\textbf{11} & 0 & 1 \\
\hline
\textbf{10} & 1 & 0 \\
\hline
\end{tabular}
\end{center}

Jedynki znajdują się więc na przekątnej mapy. W tym przypadku nie można połączyć sąsiednich jedynek w większe grupy, ponieważ żadna z nich nie posiada sąsiedniej jedynki zgodnie z zasadami map Karnaugha. Oznacza to, że bez wykorzystania dodatkowych zależności funkcja $s$ pozostaje w postaci czterech iloczynów trójzmiennych. Jednocześnie sama struktura problemu pozwala zauważyć, że funkcję można znacznie efektywniej zrealizować, rozpoznając w niej dwukrotne zastosowanie alternatywy rozłącznej (XOR).

Mapa Karnaugha dla wyjścia $c_{out}$ przyjmuje postać:

\begin{center}
\begin{tabular}{|c|c|c|}
\hline
$ab \backslash c_{in}$ & \textbf{0} & \textbf{1} \\
\hline
\textbf{00} & 0 & 0 \\
\hline
\textbf{01} & 0 & 1 \\
\hline
\textbf{11} & 1 & 1 \\
\hline
\textbf{10} & 0 & 1 \\
\hline
\end{tabular}
\end{center}

Dla funkcji $c_{out}$ możliwe jest utworzenie trzech grup po dwie jedynki: pary w wierszu $ab = 11$, pary w kolumnie $c_{in} = 1$ obejmującej wiersze $01$ i $11$ oraz pary w kolumnie $c_{in} = 1$ obejmującej wiersze $11$ i $10$. W wyniku grupowania otrzymujemy zredukowaną postać funkcji większości:
\begin{equation*}
c_{out} = a \cdot b \vee b \cdot c_{in} \vee a \cdot c_{in}.
\end{equation*}

\subsection{Minimalizacja funkcji logicznej}

Dla wyjścia sumy, na podstawie analizy struktury funkcji, otrzymujemy:
\begin{equation*}
s = a \oplus b \oplus c_{in}.
\end{equation*}

Wprowadzając sygnał pomocniczy:
\begin{equation*}
c_1 = a \oplus b,
\end{equation*}
zapis ten upraszcza się do:
\begin{equation*}
s = c_1 \oplus c_{in}.
\end{equation*}

Dla wyjścia przeniesienia, w wyniku grupowania na mapie Karnaugha:
\begin{equation*}
c_{out} = a \cdot b \vee a \cdot c_{in} \vee b \cdot c_{in}.
\end{equation*}

Alternatywnie, wykorzystując wspólny sygnał $c_1$, przeniesienie można zapisać jako:
\begin{equation*}
c_{out} = a \cdot b \vee c_1 \cdot c_{in}.
\end{equation*}

Jest to postać znacznie bardziej przejrzysta od postaci kanonicznej.

Jeżeli w realizacji można zastosować bramki XOR, każdy z członów może zostać zrealizowany pojedynczą bramką:
\begin{equation*}
c_1 = a \oplus b, \qquad s = c_1 \oplus c_{in}.
\end{equation*}

Ostateczna realizacja wymaga więc:
\begin{itemize}
    \item 2 $\times$ XOR ($c_1$ oraz $s$),
    \item 2 $\times$ AND2 ($a \cdot b$ oraz $c_1 \cdot c_{in}$),
    \item 1 $\times$ OR2 (sumowanie członów przeniesienia).
\end{itemize}

Czyli łącznie 5 bramek.

To rozwiązanie będzie naszym wariantem zoptymalizowanym pod względem liczby bramek, o ile w dalszej części przyjmiemy XOR jako dostępny element logiczny.

\subsection{Porównanie liczby bramek przed i po minimalizacji}

Dla realizacji funkcji w sposób bezpośredni, na podstawie postaci kanonicznej, konieczne byłoby zastosowanie wielu bramek AND, NOT oraz OR.

Postać kanoniczna wyjścia sumy:
\begin{equation*}
s = a' \cdot b' \cdot c_{in} \vee a' \cdot b \cdot c_{in}' \vee a \cdot b' \cdot c_{in}' \vee a \cdot b \cdot c_{in}
\end{equation*}
wymaga utworzenia czterech iloczynów trójzmiennych, a następnie ich zsumowania. Realizacja ta wymaga: 3 bramek NOT (negacje $a$, $b$, $c_{in}$), 4 bramek AND3 oraz 1 bramki OR4, czyli łącznie 8 bramek. Analogicznie dla wyjścia $c_{out}$ kolejne 8 bramek, przy czym bramki NOT mogą być wspólne.

Łącznie realizacja bezpośrednia wymaga około 13--16 bramek logicznych o głębokości trzech poziomów.

Po zastosowaniu zależności wynikającej z alternatywy rozłącznej oraz grupowania na mapach Karnaugha otrzymujemy:
\begin{equation*}
s = a \oplus b \oplus c_{in}, \qquad c_{out} = a \cdot b \vee c_1 \cdot c_{in}, \qquad c_1 = a \oplus b.
\end{equation*}

W rezultacie układ można zrealizować przy użyciu dwóch bramek XOR, dwóch bramek AND oraz jednej bramki OR, czyli łącznie 5 bramek.

Porównanie obu wariantów przedstawia tabela:

\begin{center}
\begin{tabular}{|l|c|c|}
\hline
\textbf{Wariant realizacji} & \textbf{Liczba bramek} & \textbf{Głębokość logiczna} \\
\hline
Postać kanoniczna (SOP) & ok. 13--16 & 3 poziomy \\
Wariant zminimalizowany & 5 & 3 poziomy \\
\hline
\end{tabular}
\end{center}

Zastosowanie wspólnego sygnału $c_1 = a \oplus b$ pozwala dodatkowo wyeliminować powtarzanie tej samej operacji dla obu wyjść, co stanowi dodatkową oszczędność w porównaniu z realizacją, w której każda funkcja minimalizowana jest niezależnie.

% Wymagane pakiety w dokumencie głównym:
%   \usepackage[utf8]{inputenc}  \usepackage[T1]{polski}
%   \usepackage{amsmath, amssymb, array, booktabs}
% ============================================================

\section{Opis techniczny projektu}

\subsection{Założenia projektowe}

Celem projektu jest zaprojektowanie układu kombinacyjnego realizującego funkcję sumatora pełnego, a następnie zminimalizowanie liczby wykorzystanych bramek logicznych.

Projektowany układ będzie dodawał trzy bity wejściowe: dwa bity argumentów oraz bit przeniesienia z niższego rzędu:
\begin{equation*}
a, \qquad b, \qquad c_{in},
\end{equation*}
gdzie $c_{in}$ oznacza przeniesienie generowane przez poprzedni (młodszy) stopień kaskady.

Układ posiada trzy wejścia binarne oraz dwa wyjścia. Wyjścia będą informowały o wartości bitu sumy oraz o przeniesieniu do rzędu wyższego.

Przyjęto następujące założenia:
\begin{itemize}
    \item układ dodaje trzy bity wejściowe: $a$, $b$ oraz $c_{in}$,
    \item wszystkie wejścia przyjmują wartości logiczne 0 lub 1,
    \item układ jest układem kombinacyjnym,
    \item układ posiada dwa wyjścia: bit sumy $s$ oraz przeniesienie $c_{out}$,
    \item stan wyjścia $s = 1$ oznacza, że suma wejść jest nieparzysta,
    \item stan wyjścia $c_{out} = 1$ oznacza, że $a + b + c_{in} \geq 2$,
    \item podstawowym celem projektu jest minimalizacja liczby bramek logicznych,
    \item poprawność działania zostanie sprawdzona za pomocą symulacji.
\end{itemize}

Minimalizacja układu ma na celu zmniejszenie jego złożoności. Mniejsza liczba zastosowanych bramek oznacza prostszy schemat logiczny, mniejsze zapotrzebowanie na elementy oraz potencjalnie mniejsze opóźnienia propagacji sygnału.

\subsection{Specyfikacja wejść i wyjść}

Projektowany sumator pełny posiada trzy wejścia:
\begin{itemize}
    \item $a$ -- pierwszy bit argumentu,
    \item $b$ -- drugi bit argumentu,
    \item $c_{in}$ -- przeniesienie z niższego rzędu.
\end{itemize}

Wyjściami układu są:
\begin{itemize}
    \item $s$ -- bit sumy,
    \item $c_{out}$ -- przeniesienie do rzędu wyższego.
\end{itemize}

Zależność pomiędzy wejściami i wyjściami można zapisać jako:
\begin{equation*}
s, \, c_{out} = f(a, b, c_{in}).
\end{equation*}

Przykładowo, dla:
\begin{equation*}
a = 1, \qquad b = 1, \qquad c_{in} = 0,
\end{equation*}
otrzymujemy:
\begin{equation*}
s = 0, \qquad c_{out} = 1.
\end{equation*}

Natomiast dla:
\begin{equation*}
a = 1, \qquad b = 1, \qquad c_{in} = 1,
\end{equation*}
otrzymujemy:
\begin{equation*}
s = 1, \qquad c_{out} = 1.
\end{equation*}

\vspace{0.5em}
\noindent Tabela sygnałów:

\begin{center}
\begin{tabular}{|c|c|l|}
\hline
\textbf{Sygnał} & \textbf{Rodzaj} & \textbf{Opis} \\
\hline
$a$       & wejście & pierwszy bit argumentu \\
$b$       & wejście & drugi bit argumentu \\
$c_{in}$  & wejście & przeniesienie z niższego rzędu \\
$s$       & wyjście & bit sumy \\
$c_{out}$ & wyjście & przeniesienie do rzędu wyższego \\
\hline
\end{tabular}
\end{center}

\subsection{Analiza logiczna układu}

Sumator pełny wykonuje arytmetyczne dodawanie trzech bitów, w wyniku czego może powstać wartość z zakresu od 0 do 3, wymagająca dwóch bitów zapisu: $c_{out}s$.

Aby suma była równa 1, nieparzysta liczba sygnałów wejściowych musi przyjmować wartość wysoką. Warunek ten realizuje dwuargumentowa operacja alternatywy rozłącznej, zastosowana dwukrotnie:
\begin{equation*}
s = a \oplus b \oplus c_{in}.
\end{equation*}

Przeniesienie pojawia się wtedy, gdy co najmniej dwa z trzech wejść znajdują się w stanie wysokim. Funkcję wyjściową przeniesienia można zapisać jako:
\begin{equation*}
c_{out} = a \cdot b \vee a \cdot c_{in} \vee b \cdot c_{in}.
\end{equation*}

Jest to tzw. funkcja większości, która przyjmuje wartość 1, gdy większość (co najmniej dwa) sygnałów wejściowych ma wartość 1.

Wprowadzając oznaczenie pomocnicze:
\begin{equation*}
c_1 = a \oplus b,
\end{equation*}
można obydwie funkcje zapisać w postaci:
\begin{equation*}
s = c_1 \oplus c_{in}, \qquad c_{out} = a \cdot b \vee c_1 \cdot c_{in}.
\end{equation*}

Jeżeli nie korzystamy z bramek XOR, alternatywę rozłączną można zapisać przy pomocy podstawowych operacji AND, OR i NOT:
\begin{equation*}
a \oplus b = a \cdot b' \vee a' \cdot b.
\end{equation*}

Stąd:
\begin{equation*}
s = a \cdot b' \cdot c_{in}' \vee a \cdot b \cdot c_{in} \vee a' \cdot b \cdot c_{in}' \vee a' \cdot b' \cdot c_{in}.
\end{equation*}

Ta postać będzie szczególnie istotna podczas dalszej analizy i minimalizacji układu.

\subsection{Tablica prawdy}

Ponieważ układ posiada trzy wejścia binarne, liczba możliwych kombinacji wynosi:
\begin{equation*}
2^3 = 8.
\end{equation*}

Pełna tablica prawdy przedstawia się następująco:

\begin{center}
\begin{tabular}{|c|c|c|c|c|c|}
\hline
$a$ & $b$ & $c_{in}$ & $a+b+c_{in}$ & $s$ & $c_{out}$ \\
\hline
0 & 0 & 0 & 0 & 0 & 0 \\
0 & 0 & 1 & 1 & 1 & 0 \\
0 & 1 & 0 & 1 & 1 & 0 \\
0 & 1 & 1 & 2 & 0 & 1 \\
1 & 0 & 0 & 1 & 1 & 0 \\
1 & 0 & 1 & 2 & 0 & 1 \\
1 & 1 & 0 & 2 & 0 & 1 \\
1 & 1 & 1 & 3 & 1 & 1 \\
\hline
\end{tabular}
\end{center}

Wartość wyjścia $s = 1$ występuje wyłącznie dla kombinacji:
\begin{equation*}
(0,0,1), \quad (0,1,0), \quad (1,0,0), \quad (1,1,1).
\end{equation*}

Funkcję w postaci sumy iloczynów mintermów można więc zapisać jako:
\begin{equation*}
s = \Sigma m(1, 2, 4, 7).
\end{equation*}

Po rozpisaniu:
\begin{equation*}
s = a' \cdot b' \cdot c_{in} \vee a' \cdot b \cdot c_{in}' \vee a \cdot b' \cdot c_{in}' \vee a \cdot b \cdot c_{in}.
\end{equation*}

Analogicznie dla wyjścia przeniesienia, gdzie $c_{out} = 1$ dla kombinacji $(0,1,1)$, $(1,0,1)$, $(1,1,0)$ oraz $(1,1,1)$:
\begin{equation*}
c_{out} = \Sigma m(3, 5, 6, 7),
\end{equation*}
\begin{equation*}
c_{out} = a' \cdot b \cdot c_{in} \vee a \cdot b' \cdot c_{in} \vee a \cdot b \cdot c_{in}' \vee a \cdot b \cdot c_{in}.
\end{equation*}

Jest to postać kanoniczna, która będzie punktem wyjścia do minimalizacji.

\subsection{Mapa Karnaugha}

Do uproszczenia funkcji wykorzystujemy mapy Karnaugha dla trzech zmiennych, osobno dla każdego z wyjść.

Przyjmujemy:
\begin{itemize}
    \item wiersze: pary zmiennych $ab$,
    \item kolumny: zmienna $c_{in}$,
    \item wiersze uporządkowane w kolejności kodu Graya: $00, 01, 11, 10$.
\end{itemize}

Mapa Karnaugha dla wyjścia $s$ przyjmuje postać:

\begin{center}
\begin{tabular}{|c|c|c|}
\hline
$ab \backslash c_{in}$ & \textbf{0} & \textbf{1} \\
\hline
\textbf{00} & 0 & 1 \\
\hline
\textbf{01} & 1 & 0 \\
\hline
\textbf{11} & 0 & 1 \\
\hline
\textbf{10} & 1 & 0 \\
\hline
\end{tabular}
\end{center}

Jedynki znajdują się więc na przekątnej mapy. W tym przypadku nie można połączyć sąsiednich jedynek w większe grupy, ponieważ żadna z nich nie posiada sąsiedniej jedynki zgodnie z zasadami map Karnaugha. Oznacza to, że bez wykorzystania dodatkowych zależności funkcja $s$ pozostaje w postaci czterech iloczynów trójzmiennych. Jednocześnie sama struktura problemu pozwala zauważyć, że funkcję można znacznie efektywniej zrealizować, rozpoznając w niej dwukrotne zastosowanie alternatywy rozłącznej (XOR).

Mapa Karnaugha dla wyjścia $c_{out}$ przyjmuje postać:

\begin{center}
\begin{tabular}{|c|c|c|}
\hline
$ab \backslash c_{in}$ & \textbf{0} & \textbf{1} \\
\hline
\textbf{00} & 0 & 0 \\
\hline
\textbf{01} & 0 & 1 \\
\hline
\textbf{11} & 1 & 1 \\
\hline
\textbf{10} & 0 & 1 \\
\hline
\end{tabular}
\end{center}

Dla funkcji $c_{out}$ możliwe jest utworzenie trzech grup po dwie jedynki: pary w wierszu $ab = 11$, pary w kolumnie $c_{in} = 1$ obejmującej wiersze $01$ i $11$ oraz pary w kolumnie $c_{in} = 1$ obejmującej wiersze $11$ i $10$. W wyniku grupowania otrzymujemy zredukowaną postać funkcji większości:
\begin{equation*}
c_{out} = a \cdot b \vee b \cdot c_{in} \vee a \cdot c_{in}.
\end{equation*}

\subsection{Minimalizacja funkcji logicznej}

Dla wyjścia sumy, na podstawie analizy struktury funkcji, otrzymujemy:
\begin{equation*}
s = a \oplus b \oplus c_{in}.
\end{equation*}

Wprowadzając sygnał pomocniczy:
\begin{equation*}
c_1 = a \oplus b,
\end{equation*}
zapis ten upraszcza się do:
\begin{equation*}
s = c_1 \oplus c_{in}.
\end{equation*}

Dla wyjścia przeniesienia, w wyniku grupowania na mapie Karnaugha:
\begin{equation*}
c_{out} = a \cdot b \vee a \cdot c_{in} \vee b \cdot c_{in}.
\end{equation*}

Alternatywnie, wykorzystując wspólny sygnał $c_1$, przeniesienie można zapisać jako:
\begin{equation*}
c_{out} = a \cdot b \vee c_1 \cdot c_{in}.
\end{equation*}

Jest to postać znacznie bardziej przejrzysta od postaci kanonicznej.

Jeżeli w realizacji można zastosować bramki XOR, każdy z członów może zostać zrealizowany pojedynczą bramką:
\begin{equation*}
c_1 = a \oplus b, \qquad s = c_1 \oplus c_{in}.
\end{equation*}

Ostateczna realizacja wymaga więc:
\begin{itemize}
    \item 2 $\times$ XOR ($c_1$ oraz $s$),
    \item 2 $\times$ AND2 ($a \cdot b$ oraz $c_1 \cdot c_{in}$),
    \item 1 $\times$ OR2 (sumowanie członów przeniesienia).
\end{itemize}

Czyli łącznie 5 bramek.

To rozwiązanie będzie naszym wariantem zoptymalizowanym pod względem liczby bramek, o ile w dalszej części przyjmiemy XOR jako dostępny element logiczny.

\subsection{Porównanie liczby bramek przed i po minimalizacji}

Dla realizacji funkcji w sposób bezpośredni, na podstawie postaci kanonicznej, konieczne byłoby zastosowanie wielu bramek AND, NOT oraz OR.

Postać kanoniczna wyjścia sumy:
\begin{equation*}
s = a' \cdot b' \cdot c_{in} \vee a' \cdot b \cdot c_{in}' \vee a \cdot b' \cdot c_{in}' \vee a \cdot b \cdot c_{in}
\end{equation*}
wymaga utworzenia czterech iloczynów trójzmiennych, a następnie ich zsumowania. Realizacja ta wymaga: 3 bramek NOT (negacje $a$, $b$, $c_{in}$), 4 bramek AND3 oraz 1 bramki OR4, czyli łącznie 8 bramek. Analogicznie dla wyjścia $c_{out}$ kolejne 8 bramek, przy czym bramki NOT mogą być wspólne.

Łącznie realizacja bezpośrednia wymaga około 13--16 bramek logicznych o głębokości trzech poziomów.

Po zastosowaniu zależności wynikającej z alternatywy rozłącznej oraz grupowania na mapach Karnaugha otrzymujemy:
\begin{equation*}
s = a \oplus b \oplus c_{in}, \qquad c_{out} = a \cdot b \vee c_1 \cdot c_{in}, \qquad c_1 = a \oplus b.
\end{equation*}

W rezultacie układ można zrealizować przy użyciu dwóch bramek XOR, dwóch bramek AND oraz jednej bramki OR, czyli łącznie 5 bramek.

Porównanie obu wariantów przedstawia tabela:

\begin{center}
\begin{tabular}{|l|c|c|}
\hline
\textbf{Wariant realizacji} & \textbf{Liczba bramek} & \textbf{Głębokość logiczna} \\
\hline
Postać kanoniczna (SOP) & ok. 13--16 & 3 poziomy \\
Wariant zminimalizowany & 5 & 3 poziomy \\
\hline
\end{tabular}
\end{center}

Zastosowanie wspólnego sygnału $c_1 = a \oplus b$ pozwala dodatkowo wyeliminować powtarzanie tej samej operacji dla obu wyjść, co stanowi dodatkową oszczędność w porównaniu z realizacją, w której każda funkcja minimalizowana jest niezależnie.
